% Part: normal-modal-logic
% Chapter: axioms-systems
% Section: logics-proofs

\documentclass[../../../include/open-logic-section]{subfiles}

\begin{document}

\olfileid{nml}{prf}{prf}

\olsection{\usetoken{P}{derivation} మరియు మోడల్ వ్యవస్థలు}

ముందుగా నార్మల్ మోడల్ తర్కాల కోసం \tetoken{వ్యుత్పత్తి}{!!a{derivation}}
అంటే ఏమిటో నిర్వచిస్తాం. స్థూలంగా, \tetoken{వ్యుత్పత్తి}{!!a{derivation}}
అనేది \tetoken{సూత్రాల}{!!{formula}s} క్రమం; దానిలోని ప్రతి
\tetoken{మూలకం}{!!{element}} కొన్ని \emph{స్వీకృతాలలో} ఒకటి
(లేదా దాని ప్రతిస్థాపన నిదర్శనం) అయి ఉండాలి, లేదా కొద్ది
నిగమన నియమాల్లో ఒకదాని ద్వారా మునుపటి \tetoken{మూలకాల}{!!{element}s}
నుంచి అనుసరించాలి. నార్మల్ మోడల్ తర్కాలకు సర్వసత్యాల అన్ని
నిదర్శనాలు\iftag{prvDiamond}{, \Ax{K}, మరియు~\Dual{}}{ మరియు~\Ax{K}}
స్వీకృతాలుగా పరిగణింపబడతాయి. దీని ఫలితంగా అతి చిన్న నార్మల్
మోడల్ తర్కమైన మోడల్ వ్యవస్థ~$\Log{K}$ లభిస్తుంది. అయితే ఇతర
వ్యవస్థలను పొందడానికి అదనపు స్వీకృతాలను చేర్చుకోవచ్చు.
నియమాలు ఎప్పుడూ మోడస్ పోనెన్స్~\MP{} మరియు అవశ్యకీకరణ~\Nec.

\begin{defn}
  మోడల్ వ్యవస్థ $\Log{K} !A_1 \dots !A_n$, అలాగే
  \tetoken{ఒక సూత్రం}{!!a{formula}}~$!B$ ఇచ్చినప్పుడు,
  $!B$ అనేది $\Log{K} !A_1 \dots !A_n$లో
  \emph{\tetoken{వ్యుత్పాదించదగినది}{!!{derivable}}} అంటాం;
  సంకేతంగా $\Log{K} !A_1 \dots !A_n \Proves !B$. ఇది అప్పుడూ అప్పుడే:
  $!C_k = !B$ అయ్యే $!C_1$, \dots, $!C_k$ అనే
  \tetoken{సూత్రాలు}{!!{formula}s} ఉండి, ప్రతి $!C_i$
  సర్వసత్య ప్రతిస్థాపన నిదర్శనం అయి ఉండాలి, లేదా
  $\Ax{K}$,\iftag{prvDiamond}{ $\Dual$,}{} $!A_1$,
  \dots, $!A_n$లలో ఒకదాని నిదర్శనం అయి ఉండాలి, లేదా
  మునుపటి \tetoken{సూత్రాల}{!!{formula}s} నుంచి \MP{} లేదా~\Nec{}
  నియమం ద్వారా అనుసరించాలి.
\end{defn}

కింది ప్రతిపాదన వల్ల $!B$ యొక్క $\Sigma$-\tetoken{వ్యుత్పత్తి}{!!{derivation}}
చూపించి $!B \in \Sigma$ అని చూపవచ్చు.

\begin{prop}
  $\Log{K} !A_1 \dots !A_n = \Setabs{!B}{\Log{K} !A_1 \dots !A_n \Proves !B}$.
\end{prop}

\begin{proof}
  \tetoken{వ్యుత్పత్తుల}{!!{derivation}s} పొడవుపై ఆగమనం చేసి
  $\Setabs{!B}{\Log{K} !A_1 \dots !A_n \Proves !B} \subseteq
  \Log{K} !A_1 \dots !A_n$ అని చూపుతాం.

  $!B$ యొక్క \tetoken{వ్యుత్పత్తి}{!!{derivation}} పొడవు~$1$
  అయితే, అందులో ఒకే \tetoken{సూత్రం}{!!{formula}} ఉంటుంది.
  ఆ \tetoken{సూత్రం}{!!{formula}} మునుపటి సూత్రాల నుంచి
  \MP{} లేదా \Nec{} ద్వారా అనుసరించలేదు; కనుక అది సర్వసత్య
  ప్రతిస్థాపన నిదర్శనం, \Ax{K}\iftag{prvDiamond}{ లేదా $\Dual$}{},
  లేదా $!A_1$, \dots,~$!A_n$లో ఒకదాని నిదర్శనం అయి ఉండాలి.
  $\Log{K}!A_1\dots!A_n$లో ఇవన్నీ ఉన్నాయి కనుక,
  $!B \in \Log{K}!A_1 \dots !A_n$.

  $!B$ యొక్క \tetoken{వ్యుత్పత్తి}{!!{derivation}} పొడవు~$> 1$
  అయితే, \tetoken{వ్యుత్పత్తిలో}{!!{derivation}} చివరి పంక్తి కాని
  \tetoken{సూత్రాల}{!!{formula}s} నుంచి
  \MP{} లేదా \Nec{} ద్వారా కూడా $!B$ రావచ్చు. $!B$, $!C$
  మరియు $!C \lif !B$ నుంచి (\MP{} ద్వారా) అనుసరిస్తే,
  ఆగమన పరికల్పన వల్ల $!C$, $!C \lif !B \in
  \Log{K}!A_1 \dots !A_n$. ప్రతి మోడల్ తర్కం మోడస్ పోనెన్స్
  కింద సంవృతం కనుక $!B \in \Log{K}!A_1 \dots
  !A_n$. $!B \equiv \Box !C$ అనేది $!C$ నుంచి \Nec{} ద్వారా
  అనుసరిస్తే, ఆగమన పరికల్పన వల్ల $!C
  \in \Log{K}!A_1 \dots !A_n$. ప్రతి నార్మల్ మోడల్ తర్కం
  $\Nec$ కింద సంవృతం కనుక, $!B \in
  \Log{K}!A_1\dots!A_n$.

  విలోమ సమితి చేరిక కోసం
  $\Sigma = \Setabs{!B}{\Log{K} !A_1 \dots !A_n \Proves !B }$
  అనేది $!A_1$, \dots, $!A_n$ల అన్ని నిదర్శనాలను కలిగి ఉండే
  నార్మల్ మోడల్ తర్కమని చూపుతాం. నిర్వచనం ప్రకారం
  $\Log{K} !A_1 \dots !A_n$ అటువంటి తర్కాల్లో అతి చిన్నది.
  \begin{enumerate}
    \item ప్రతి సర్వసత్యం~$!B$ సర్వసత్య ప్రతిస్థాపన నిదర్శనమే;
      కనుక $\Log{K}!A_1\dots!A_n \Proves !B$. అందువల్ల $\Sigma$
      అన్ని సర్వసత్యాలను కలిగి ఉంటుంది.
    \item $\Log{K}!A_1\dots!A_n \Proves !C$ మరియు
      $\Log{K}!A_1\dots!A_n \Proves !C \lif !B$ అయితే,
      $\Log{K}!A_1\dots!A_n \Proves !B$: $!C$ యొక్క
      \tetoken{వ్యుత్పత్తిని}{!!{derivation}}, $!C \lif !B$ యొక్క
      వ్యుత్పత్తితో కలిపి, చివరగా~$!B$ పంక్తిని చేర్చండి.
      ఆ చివరి పంక్తికి \MP{} సమర్థన. అందువల్ల $\Sigma$
      మోడస్ పోనెన్స్ కింద సంవృతం.
    \item $!B$కు \tetoken{ఒక వ్యుత్పత్తి}{!!a{derivation}} ఉంటే,
      $!B$ యొక్క ప్రతి ప్రతిస్థాపన నిదర్శనానికీ
      \tetoken{ఒక వ్యుత్పత్తి}{!!a{derivation}} ఉంటుంది:
      ఆ ప్రతిస్థాపనను \tetoken{వ్యుత్పత్తిలోని}{!!{derivation}}
      ప్రతి \tetoken{సూత్రానికీ}{!!{formula}} వర్తింపజేయండి.
      (వ్యాయామం: \tetoken{వ్యుత్పత్తుల}{!!{derivation}s} పొడవుపై
      ఆగమనం చేసి ఫలితం కూడా సరైన
      \tetoken{వ్యుత్పత్తి}{!!{derivation}} అని నిరూపించండి.)
      అందువల్ల $\Sigma$ ఏకరీతి ప్రతిస్థాపన కింద సంవృతం.
      (ఇప్పటికి $\Sigma$ మోడల్ తర్కపు అన్ని షరతులనూ
      నెరవేరుస్తుందని స్థాపించాం.)
    \item $\Log{K}!A_1\dots!A_n \Proves \Ax{K}$;
      కనుక $\Ax{K} \in \Sigma$.
      \sourcecorrection{OLTENMLAXSPRF-001}{స్థిర మూలంలో ఈ
        సభ్యత్వ ప్రకటనలో స్వీకృత పథకం పేరును సూత్రంలా వాడింది.
        వ్యుత్పాదించబడిన స్వీకృత సూత్రాన్నే సమితి సభ్యునిగా
        ఇక్కడ స్పష్టంగా రాశాం.}
    \tagitem{prvDiamond}{మనకు
      $\Log{K}!A_1\dots!A_n \Proves \Dual$; కనుక $\Dual \in \Sigma$.}{}
    \item $\Log{K}!A_1\dots!A_n \Proves !C$ అయితే,
      అదనపు పంక్తి~$\Box !C$కు~\Nec{} సమర్థన ఉంటుంది.
      అందువల్ల $\Sigma$,~\Nec{} కింద సంవృతం. కనుక
      $\Sigma$ నార్మల్.
  \end{enumerate}
\end{proof}

\end{document}
